\cite{gupta} states that the CLV is generally defined as the present value of all future profits obtained from a customer over his or her life of relationship with a firm. With CLV for a customer defined as \citep{gupta}

\[
CLV = \sum_{t=0}^T \frac{(p_t - c_t)r_t}{(1+i)^t} - AC
\]

\noindent where
\begin{align*}
    &p_t = \text{price paid by a consumer at time t,} \\
    &c_t = \text{direct cost of servicing the customer at time t,} \\
    &i = \text{discount rate or cost of capital for the firm,} \\
    &r_t = \text{probability of customer repeat buying or being
“alive” at time t,} \\
    &AC = \text{acquisition cost, and} \\
    &T = \text{time horizon for estimating CLV} \\
\end{align*}